% !TEX root = ../main.tex
% =====================================================================
% 10. LIMITATIONS AND CONCLUSION -- page budget 0.4
%
% Limitations come FIRST and are stated flatly. A reviewer who finds a
% limitation we already named reads it as honesty; one who finds a
% limitation we hid stops trusting the rest.
% =====================================================================

\section{Limitations}
\label{sec:limitations}

Five bounds apply to everything above. \textbf{Explanation correctness is measured on one synthetic
dataset}, because \ndcg{} needs per-step ground truth and only a generator supplies it; how these
results hold on real, noisier signals is untested, and constructing a second ground-truth dataset is
the single most valuable extension we can name. \textbf{Finding~\ref{find:selection} rests on
single-seed runs}: with a per-model spread of $\pm0.016$ accuracy we read only the large steps in
\Cref{tab:kappa} and refuse the small ones, but repeated runs would settle the shape of the curve
properly. \textbf{The headline models were selected on the dataset they are headlined on}
(\Cref{sec:protocol}), which is exactly what \Cref{tab:ucr} exists to bound. \textbf{Only univariate
series were tested}; the encoder is indifferent to input channel count, so multivariate input should
work mechanically, but that was not verified. \textbf{Faithfulness is not usefulness}: AOPCR and
\ndcg{} both ask whether the stated explanation matches something checkable, and neither establishes
that a person would find the highlighted steps informative -- that requires a human study.

\section{Conclusion}
\label{sec:conclusion}

We treated a family of inherently interpretable time series classifiers as a design space rather
than as a sequence of models, and swept it under one configuration. Three of the four results point
the same way: the aggregator matters more than its reputation suggests. Selection inside the
aggregation rule, not encoder capacity, is what makes an explanation rank the right time steps, and
it costs about one accuracy point (\Cref{sec:selection}). The head axis moves accuracy as much as the
encoder axis, and the two interact, so the pair has to be searched (\Cref{sec:cheap}). And a head's
reduction property protects it only if the optimiser is given a reason to use it
(\Cref{sec:safety}).

The fourth result is a caution about how this work is measured. AOPCR moves by $5\times$ under a
change of training recipe alone (\Cref{sec:aopcr}), which is larger than any difference between the
methods we compared. Until a normalised replacement exists, faithfulness claims in this literature
should be made within a single training configuration, against baselines re-trained inside it. That
is a heavier experimental burden than quoting a published number, and \Cref{tab:aopcr} is our
argument that it is not optional.

Finally, interpretability turned out to be cheap and capacity turned out not to be: a $41$\,K model
matched a $424$\,K baseline on accuracy and both explanation metrics simultaneously, while the same
architecture widened back up recovered the archive-wide performance the small version lost. The gap
that remains against the best result we measured is not an architectural one at all -- it is a
training budget of $1500$ epochs against $400$.
